% Options for packages loaded elsewhere
\PassOptionsToPackage{unicode}{hyperref}
\PassOptionsToPackage{hyphens}{url}
\PassOptionsToPackage{dvipsnames,svgnames,x11names}{xcolor}
\documentclass[
  11pt,
  a4paper,
]{article}
\usepackage{xcolor}
\usepackage[margin=1in]{geometry}
\usepackage{amsmath,amssymb}
\setcounter{secnumdepth}{-\maxdimen} % remove section numbering
\usepackage{iftex}
\ifPDFTeX
  \usepackage[T1]{fontenc}
  \usepackage[utf8]{inputenc}
  \usepackage{textcomp} % provide euro and other symbols
\else % if luatex or xetex
  \usepackage{unicode-math} % this also loads fontspec
  \defaultfontfeatures{Scale=MatchLowercase}
  \defaultfontfeatures[\rmfamily]{Ligatures=TeX,Scale=1}
\fi
\usepackage{lmodern}
\ifPDFTeX\else
  % xetex/luatex font selection
  \ifLuaTeX
    \usepackage{luaotfload}
    \directlua{luaotfload.add_fallback("mainfontfallback",{
      "NotoSerifBengali:mode=harf"
    })}
  \fi
  \setmainfont[,RawFeature={fallback=mainfontfallback}]{Latin Modern Roman}
  \ifLuaTeX
    \usepackage{luaotfload}
    \directlua{luaotfload.add_fallback("sansfontfallback",{
      "NotoSansBengali:mode=harf"
    })}
  \fi
  \setsansfont[,RawFeature={fallback=sansfontfallback}]{Latin Modern Sans}
  \ifLuaTeX
    \usepackage{luaotfload}
    \directlua{luaotfload.add_fallback("monofontfallback",{
      "NotoSansBengali:mode=harf"
    })}
  \fi
  \setmonofont[,RawFeature={fallback=monofontfallback}]{DejaVu Sans Mono}
\fi
% Use upquote if available, for straight quotes in verbatim environments
\IfFileExists{upquote.sty}{\usepackage{upquote}}{}
\IfFileExists{microtype.sty}{% use microtype if available
  \usepackage[]{microtype}
  \UseMicrotypeSet[protrusion]{basicmath} % disable protrusion for tt fonts
}{}
\makeatletter
\@ifundefined{KOMAClassName}{% if non-KOMA class
  \IfFileExists{parskip.sty}{%
    \usepackage{parskip}
  }{% else
    \setlength{\parindent}{0pt}
    \setlength{\parskip}{6pt plus 2pt minus 1pt}}
}{% if KOMA class
  \KOMAoptions{parskip=half}}
\makeatother
\usepackage{color}
\usepackage{fancyvrb}
\newcommand{\VerbBar}{|}
\newcommand{\VERB}{\Verb[commandchars=\\\{\}]}
\DefineVerbatimEnvironment{Highlighting}{Verbatim}{commandchars=\\\{\}}
% Add ',fontsize=\small' for more characters per line
\usepackage{framed}
\definecolor{shadecolor}{RGB}{248,248,248}
\newenvironment{Shaded}{\begin{snugshade}}{\end{snugshade}}
\newcommand{\AlertTok}[1]{\textcolor[rgb]{0.94,0.16,0.16}{#1}}
\newcommand{\AnnotationTok}[1]{\textcolor[rgb]{0.56,0.35,0.01}{\textbf{\textit{#1}}}}
\newcommand{\AttributeTok}[1]{\textcolor[rgb]{0.13,0.29,0.53}{#1}}
\newcommand{\BaseNTok}[1]{\textcolor[rgb]{0.00,0.00,0.81}{#1}}
\newcommand{\BuiltInTok}[1]{#1}
\newcommand{\CharTok}[1]{\textcolor[rgb]{0.31,0.60,0.02}{#1}}
\newcommand{\CommentTok}[1]{\textcolor[rgb]{0.56,0.35,0.01}{\textit{#1}}}
\newcommand{\CommentVarTok}[1]{\textcolor[rgb]{0.56,0.35,0.01}{\textbf{\textit{#1}}}}
\newcommand{\ConstantTok}[1]{\textcolor[rgb]{0.56,0.35,0.01}{#1}}
\newcommand{\ControlFlowTok}[1]{\textcolor[rgb]{0.13,0.29,0.53}{\textbf{#1}}}
\newcommand{\DataTypeTok}[1]{\textcolor[rgb]{0.13,0.29,0.53}{#1}}
\newcommand{\DecValTok}[1]{\textcolor[rgb]{0.00,0.00,0.81}{#1}}
\newcommand{\DocumentationTok}[1]{\textcolor[rgb]{0.56,0.35,0.01}{\textbf{\textit{#1}}}}
\newcommand{\ErrorTok}[1]{\textcolor[rgb]{0.64,0.00,0.00}{\textbf{#1}}}
\newcommand{\ExtensionTok}[1]{#1}
\newcommand{\FloatTok}[1]{\textcolor[rgb]{0.00,0.00,0.81}{#1}}
\newcommand{\FunctionTok}[1]{\textcolor[rgb]{0.13,0.29,0.53}{\textbf{#1}}}
\newcommand{\ImportTok}[1]{#1}
\newcommand{\InformationTok}[1]{\textcolor[rgb]{0.56,0.35,0.01}{\textbf{\textit{#1}}}}
\newcommand{\KeywordTok}[1]{\textcolor[rgb]{0.13,0.29,0.53}{\textbf{#1}}}
\newcommand{\NormalTok}[1]{#1}
\newcommand{\OperatorTok}[1]{\textcolor[rgb]{0.81,0.36,0.00}{\textbf{#1}}}
\newcommand{\OtherTok}[1]{\textcolor[rgb]{0.56,0.35,0.01}{#1}}
\newcommand{\PreprocessorTok}[1]{\textcolor[rgb]{0.56,0.35,0.01}{\textit{#1}}}
\newcommand{\RegionMarkerTok}[1]{#1}
\newcommand{\SpecialCharTok}[1]{\textcolor[rgb]{0.81,0.36,0.00}{\textbf{#1}}}
\newcommand{\SpecialStringTok}[1]{\textcolor[rgb]{0.31,0.60,0.02}{#1}}
\newcommand{\StringTok}[1]{\textcolor[rgb]{0.31,0.60,0.02}{#1}}
\newcommand{\VariableTok}[1]{\textcolor[rgb]{0.00,0.00,0.00}{#1}}
\newcommand{\VerbatimStringTok}[1]{\textcolor[rgb]{0.31,0.60,0.02}{#1}}
\newcommand{\WarningTok}[1]{\textcolor[rgb]{0.56,0.35,0.01}{\textbf{\textit{#1}}}}
\usepackage{longtable,booktabs,array}
\usepackage{calc} % for calculating minipage widths
% Correct order of tables after \paragraph or \subparagraph
\usepackage{etoolbox}
\makeatletter
\patchcmd\longtable{\par}{\if@noskipsec\mbox{}\fi\par}{}{}
\makeatother
% Allow footnotes in longtable head/foot
\IfFileExists{footnotehyper.sty}{\usepackage{footnotehyper}}{\usepackage{footnote}}
\makesavenoteenv{longtable}
\setlength{\emergencystretch}{3em} % prevent overfull lines
\providecommand{\tightlist}{%
  \setlength{\itemsep}{0pt}\setlength{\parskip}{0pt}}
% paper_preamble.tex — shared LaTeX preamble for the PCMA/SI paper set.
%
% Compiled with lualatex (NOT pdflatex) because several papers contain
% Bengali script: the CLL paper realises one cognitive state into
% English, Bengali and Banglish, so the Bengali output is EVIDENCE and
% cannot be transliterated away without destroying the point being made.
%
% Font selection is NOT done here. pandoc >= 3 has native
% `mainfontfallback` / `monofontfallback` variables that emit the
% luaotfload fallback registration in the correct position relative to
% unicode-math; doing it by hand in this header lands after unicode-math
% and breaks font loading outright. See tools/build_papers.sh.

\usepackage{amsmath}
\usepackage{amssymb}
\usepackage{mathtools}   % \xleftrightarrow etc. used in cll.md §5
\usepackage{booktabs}

% ---------------------------------------------------------------------
% Unicode symbols appearing in PROSE rather than inside $...$.
%
% The papers were written as markdown and read as markdown, so they use
% real characters — "the ω-boost", "α-child", "ΔR coherence" — instead of
% LaTeX macros. Under unicode-math those codepoints are math-mode-only
% and abort the build ("\symbballowed only in math mode"). Mapping each
% to \ensuremath keeps the markdown readable on GitHub AND compilable,
% which is the whole point of keeping one source.
% ---------------------------------------------------------------------
\usepackage{newunicodechar}

% Greek, lower case
\newunicodechar{α}{\ensuremath{\alpha}}
\newunicodechar{β}{\ensuremath{\beta}}
\newunicodechar{γ}{\ensuremath{\gamma}}
\newunicodechar{δ}{\ensuremath{\delta}}
\newunicodechar{ε}{\ensuremath{\varepsilon}}
\newunicodechar{η}{\ensuremath{\eta}}
\newunicodechar{θ}{\ensuremath{\theta}}
\newunicodechar{κ}{\ensuremath{\kappa}}
\newunicodechar{λ}{\ensuremath{\lambda}}
\newunicodechar{μ}{\ensuremath{\mu}}
\newunicodechar{π}{\ensuremath{\pi}}
\newunicodechar{ρ}{\ensuremath{\rho}}
\newunicodechar{σ}{\ensuremath{\sigma}}
\newunicodechar{τ}{\ensuremath{\tau}}
\newunicodechar{φ}{\ensuremath{\varphi}}
\newunicodechar{χ}{\ensuremath{\chi}}
\newunicodechar{ψ}{\ensuremath{\psi}}
\newunicodechar{ω}{\ensuremath{\omega}}
\newunicodechar{ν}{\ensuremath{\nu}}
\newunicodechar{ξ}{\ensuremath{\xi}}
\newunicodechar{ζ}{\ensuremath{\zeta}}
\newunicodechar{ι}{\ensuremath{\iota}}

% Greek, upper case
\newunicodechar{Δ}{\ensuremath{\Delta}}
\newunicodechar{Φ}{\ensuremath{\Phi}}
\newunicodechar{Σ}{\ensuremath{\Sigma}}
\newunicodechar{Ψ}{\ensuremath{\Psi}}
\newunicodechar{Ω}{\ensuremath{\Omega}}
\newunicodechar{Γ}{\ensuremath{\Gamma}}
\newunicodechar{Λ}{\ensuremath{\Lambda}}
\newunicodechar{Θ}{\ensuremath{\Theta}}
\newunicodechar{Π}{\ensuremath{\Pi}}

% Operators and relations
\newunicodechar{×}{\ensuremath{\times}}
\newunicodechar{·}{\ensuremath{\cdot}}
\newunicodechar{±}{\ensuremath{\pm}}
\newunicodechar{−}{\ensuremath{-}}
\newunicodechar{≈}{\ensuremath{\approx}}
\newunicodechar{≤}{\ensuremath{\leq}}
\newunicodechar{≥}{\ensuremath{\geq}}
\newunicodechar{≠}{\ensuremath{\neq}}
\newunicodechar{∈}{\ensuremath{\in}}
\newunicodechar{∂}{\ensuremath{\partial}}
\newunicodechar{∇}{\ensuremath{\nabla}}
\newunicodechar{√}{\ensuremath{\sqrt{}}}
\newunicodechar{∞}{\ensuremath{\infty}}
\newunicodechar{∝}{\ensuremath{\propto}}
\newunicodechar{∑}{\ensuremath{\sum}}
\newunicodechar{∫}{\ensuremath{\int}}
\newunicodechar{ℝ}{\ensuremath{\mathbb{R}}}
\newunicodechar{ℕ}{\ensuremath{\mathbb{N}}}
\newunicodechar{‖}{\ensuremath{\|}}
\newunicodechar{⊕}{\ensuremath{\oplus}}
\newunicodechar{⊗}{\ensuremath{\otimes}}
\newunicodechar{∀}{\ensuremath{\forall}}
\newunicodechar{∃}{\ensuremath{\exists}}
\newunicodechar{¬}{\ensuremath{\neg}}
\newunicodechar{∧}{\ensuremath{\wedge}}
\newunicodechar{∨}{\ensuremath{\vee}}

% Arrows
\newunicodechar{→}{\ensuremath{\rightarrow}}
\newunicodechar{←}{\ensuremath{\leftarrow}}
\newunicodechar{↔}{\ensuremath{\leftrightarrow}}
\newunicodechar{↑}{\ensuremath{\uparrow}}
\newunicodechar{↓}{\ensuremath{\downarrow}}
\newunicodechar{⇒}{\ensuremath{\Rightarrow}}
\newunicodechar{⇐}{\ensuremath{\Leftarrow}}
\newunicodechar{⇔}{\ensuremath{\Leftrightarrow}}

% Super/subscripts written as literal characters
\newunicodechar{²}{\ensuremath{^{2}}}
\newunicodechar{³}{\ensuremath{^{3}}}
\newunicodechar{⁴}{\ensuremath{^{4}}}
\newunicodechar{⁵}{\ensuremath{^{5}}}
\newunicodechar{⁶}{\ensuremath{^{6}}}
\newunicodechar{₀}{\ensuremath{_{0}}}
\newunicodechar{₁}{\ensuremath{_{1}}}
\newunicodechar{₂}{\ensuremath{_{2}}}
\newunicodechar{₃}{\ensuremath{_{3}}}
\newunicodechar{ᵀ}{\ensuremath{^{\mathsf{T}}}}

% Status glyphs. The papers use emoji as status markers in tables
% (implemented / failed / warning). Mapped to typographic equivalents:
% a colour emoji font in lualatex is unreliable, and the MEANING is what
% carries — a reader needs to see "this one failed", not a red cross.
\newunicodechar{✅}{\ensuremath{\checkmark}}
\newunicodechar{✓}{\ensuremath{\checkmark}}
\newunicodechar{❌}{\ensuremath{\times}}
\newunicodechar{✗}{\ensuremath{\times}}
\newunicodechar{⚠}{\textbf{!}}
\newunicodechar{🧩}{\textbf{[piece]}}
\newunicodechar{📋}{\textbf{[list]}}
\newunicodechar{🟡}{\textbf{[partial]}}
\newunicodechar{🔴}{\textbf{[dead]}}
\newunicodechar{🟢}{\textbf{[live]}}
\newunicodechar{▼}{\ensuremath{\blacktriangledown}}
\newunicodechar{►}{\ensuremath{\blacktriangleright}}
\newunicodechar{□}{\ensuremath{\square}}
\newunicodechar{█}{\rule{0.6em}{0.9ex}}
\newunicodechar{•}{\textbullet}
\newunicodechar{†}{\textdagger}
\newunicodechar{‑}{-}
\newunicodechar{‐}{-}
\newunicodechar{…}{\ldots}
\newunicodechar{।}{\textbar}   % danda: rendered by the Bengali fallback in text
\newunicodechar{⁻}{\ensuremath{^{-}}}

% Zero-width formatting codepoints. These are invisible in an editor but
% are real characters to LaTeX: an emoji written as "⚠️" is TWO
% codepoints (U+26A0 U+FE0F), so declaring only the visible one leaves a
% variation selector with no definition and the build dies on a symbol
% nobody can see. Both are mapped to nothing.
\newunicodechar{^^^^fe0f}{}    % VARIATION SELECTOR-16
\newunicodechar{^^^^200d}{}    % ZERO WIDTH JOINER
\newunicodechar{^^^^200b}{}    % ZERO WIDTH SPACE

\usepackage{longtable}
\usepackage{array}

% Comparison matrices and cost ledgers are wide. Let them shrink to the
% text block rather than overflow the margin.
\usepackage{adjustbox}

% Code blocks: wrap long lines instead of running off the page. Several
% papers quote real C++ with long identifiers.
\usepackage{fvextra}
\fvset{breaklines=true,breakanywhere=true,fontsize=\small}

\setlength{\parskip}{0.5em}
\setlength{\parindent}{0pt}

\usepackage{fancyhdr}
\pagestyle{fancy}
\fancyhf{}
\fancyfoot[C]{\thepage}
\fancyhead[L]{\small\itshape Synthetic Intelligence — research lineage}
\fancyhead[R]{\small\itshape irfan@furylogic.com}
\renewcommand{\headrulewidth}{0.2pt}
\setlength{\headheight}{14pt}
\setlength{\headsep}{18pt}
\usepackage{bookmark}
\IfFileExists{xurl.sty}{\usepackage{xurl}}{} % add URL line breaks if available
\urlstyle{same}
\hypersetup{
  pdftitle={The ZeroGrad Cognitive Engine: A Backprop-Free, Non-Numerical Training Paradigm for Physics-Based Reasoning},
  pdfauthor={Irfan Mahir --- Independent Researcher, Dhaka, Bangladesh --- irfan@furylogic.com},
  colorlinks=true,
  linkcolor={black},
  filecolor={Maroon},
  citecolor={Blue},
  urlcolor={[HTML]{0B5394}},
  pdfcreator={LaTeX via pandoc}}

\title{The ZeroGrad Cognitive Engine: A Backprop-Free, Non-Numerical
Training Paradigm for Physics-Based Reasoning}
\author{Irfan Mahir --- Independent Researcher, Dhaka, Bangladesh ---
irfan@furylogic.com}
\date{}

\begin{document}
\maketitle

{
\hypersetup{linkcolor=}
\setcounter{tocdepth}{3}
\tableofcontents
}
\section{The ZeroGrad Cognitive Engine: A Backprop-Free, Non-Numerical
Training Paradigm for Physics-Based
Reasoning}\label{the-zerograd-cognitive-engine-a-backprop-free-non-numerical-training-paradigm-for-physics-based-reasoning}

\textbf{Irfan Mahir}\\
\emph{Independent Researcher, Dhaka, Bangladesh}\\
irfan@furylogic.com

\begin{quote}
\emph{``Thinking and understanding is the art of innovation and
invention.''}
\end{quote}

\begin{center}\rule{0.5\linewidth}{0.5pt}\end{center}

\subsection{Abstract}\label{abstract}

We present the \textbf{ZeroGrad Cognitive Engine (ZCE)} , a novel
training paradigm for the Predictive Cognitive Manifold Architecture
(P-CMA) and Cognitive Continuum Hypothesis (CCH) that completely
eliminates the need for both backpropagation and numerical gradient
computation. ZCE replaces global gradient-based optimization with a
\textbf{local, forward-only learning framework} where each
component---Encoder, Transition, and Potential---is trained
independently using local objectives and Hebbian-style update rules.

Unlike Transformers that require backpropagation through billions of
parameters, and unlike numerical gradient methods that scale as
\(O(N_{\text{params}} \times N_{\text{network}})\), ZCE scales as
\textbf{\(O(N_{\text{network}})\)} ---enabling training of 130M+
parameter models on consumer hardware without the memory or time
penalties of traditional methods.

We demonstrate that ZCE preserves all key properties of P-CMA and CCH: -
\textbf{Non-autoregressive inference} (CEM planning, not greedy
generation) - \textbf{O(1) memory} (state compression, no KV-cache) -
\textbf{100\% logical transitivity} (guaranteed by
\(\nabla \times F = 0\)) - \textbf{No hallucination} (energy barrier
prevents inconsistency)

The ZeroGrad paradigm represents a fundamental departure from the
gradient-centric training of deep learning, offering a path toward truly
scalable, interpretable, and biologically plausible cognitive
architectures.

\begin{center}\rule{0.5\linewidth}{0.5pt}\end{center}

\subsection{1. Introduction}\label{introduction}

\subsubsection{1.1 The Gradient
Bottleneck}\label{the-gradient-bottleneck}

Modern AI systems are built on two pillars: \textbf{backpropagation} and
\textbf{numerical optimization}. Both are fundamentally limited:

\begin{longtable}[]{@{}
  >{\raggedright\arraybackslash}p{(\linewidth - 4\tabcolsep) * \real{0.3333}}
  >{\raggedright\arraybackslash}p{(\linewidth - 4\tabcolsep) * \real{0.3333}}
  >{\raggedright\arraybackslash}p{(\linewidth - 4\tabcolsep) * \real{0.3333}}@{}}
\toprule\noalign{}
\begin{minipage}[b]{\linewidth}\raggedright
Method
\end{minipage} & \begin{minipage}[b]{\linewidth}\raggedright
Scaling
\end{minipage} & \begin{minipage}[b]{\linewidth}\raggedright
Problem
\end{minipage} \\
\midrule\noalign{}
\endhead
\bottomrule\noalign{}
\endlastfoot
\textbf{Backpropagation} & \(O(N_{\text{network}})\) & Requires storing
all activations; introduces vanishing/exploding gradients; couples all
layers \\
\textbf{Numerical Gradients} &
\(O(N_{\text{params}} \times N_{\text{network}})\) & \textbf{Impossible
for large models} (38 days per step for 130M params) \\
\end{longtable}

\textbf{The core insight:} Both methods require global, coordinated
updates across all parameters. This is inefficient, brittle, and
biologically implausible.

\subsubsection{1.2 The Local Learning
Alternative}\label{the-local-learning-alternative}

Recent research on the \textbf{Forward-Forward Algorithm} and
\textbf{Local Learning Rules} has shown that neural networks can be
trained without backpropagation. Each layer learns to maximize a local
``goodness'' function, using only forward passes.

We extend this idea to the P-CMA architecture, where each component has
a natural local objective:

\begin{longtable}[]{@{}
  >{\raggedright\arraybackslash}p{(\linewidth - 2\tabcolsep) * \real{0.5000}}
  >{\raggedright\arraybackslash}p{(\linewidth - 2\tabcolsep) * \real{0.5000}}@{}}
\toprule\noalign{}
\begin{minipage}[b]{\linewidth}\raggedright
Component
\end{minipage} & \begin{minipage}[b]{\linewidth}\raggedright
Local Objective
\end{minipage} \\
\midrule\noalign{}
\endhead
\bottomrule\noalign{}
\endlastfoot
\textbf{Encoder} & Maximize goodness (low energy) for consistent
inputs \\
\textbf{Transition T(c,a)} & Minimize local prediction error
\(\|c_{next} - T(c,a)\|^2\) \\
\textbf{Potential Φ(c)} & Minimize energy for consistent states,
maximize for inconsistent states \\
\end{longtable}

\textbf{The ZeroGrad paradigm trains these components independently,
using local, forward-only updates.}

\begin{center}\rule{0.5\linewidth}{0.5pt}\end{center}

\subsection{2. Mathematical Formulation}\label{mathematical-formulation}

\subsubsection{2.1 The ZeroGrad Principle}\label{the-zerograd-principle}

\textbf{Definition 2.1 (ZeroGrad Training).} ZeroGrad training is a
backprop-free, numerical-gradient-free learning framework where each
parameter \(\theta_i\) is updated based on a local objective \(L_i\)
that depends only on the activity of the layer itself:

\[\Delta \theta_i = \eta \cdot \frac{\partial L_i}{\partial \theta_i}\]

\textbf{Crucially:} The gradient of \(L_i\) is computed
\textbf{analytically} for that layer alone, without backpropagating
through the rest of the network. This is \(O(1)\) per layer, not
\(O(N_{\text{params}})\).

\subsubsection{2.2 The Local Goodness
Function}\label{the-local-goodness-function}

For each layer, we define a \textbf{goodness function} \(G_i\) that is
maximized during training:

\begin{longtable}[]{@{}
  >{\raggedright\arraybackslash}p{(\linewidth - 4\tabcolsep) * \real{0.3333}}
  >{\raggedright\arraybackslash}p{(\linewidth - 4\tabcolsep) * \real{0.3333}}
  >{\raggedright\arraybackslash}p{(\linewidth - 4\tabcolsep) * \real{0.3333}}@{}}
\toprule\noalign{}
\begin{minipage}[b]{\linewidth}\raggedright
Component
\end{minipage} & \begin{minipage}[b]{\linewidth}\raggedright
Goodness Function
\end{minipage} & \begin{minipage}[b]{\linewidth}\raggedright
Target
\end{minipage} \\
\midrule\noalign{}
\endhead
\bottomrule\noalign{}
\endlastfoot
\textbf{Encoder} & \(G_E = -\Phi(\text{Encoder}(x))\) & Maximize for
consistent inputs \\
\textbf{Transition} & \(G_T = -\|c_{next} - T(c,a)\|^2\) & Minimize
prediction error \\
\textbf{Potential} & \(G_\Phi = -\Phi(c)\) & Minimize energy for
consistent states \\
\end{longtable}

\textbf{The update rule for any layer is:}

\[\Delta \theta_i = \eta \cdot \frac{\partial G_i}{\partial \theta_i}\]

This requires only the local gradient of \(G_i\) with respect to
\(\theta_i\), which is computed analytically.

\subsubsection{2.3 The Forward-Forward Rule for Encoder and
Potential}\label{the-forward-forward-rule-for-encoder-and-potential}

Following the Forward-Forward algorithm, we train the Encoder using
\textbf{positive and negative passes}:

\textbf{Positive Pass:} Sample a consistent input \(x^+\) (e.g., a valid
premise). Compute \(c^+ = \text{Encoder}(x^+)\). Update weights to
\textbf{maximize} \(G_E = -\Phi(c^+)\).

\textbf{Negative Pass:} Sample an inconsistent input \(x^-\) (e.g., a
random/corrupted premise). Compute \(c^- = \text{Encoder}(x^-)\). Update
weights to \textbf{minimize} \(G_E = -\Phi(c^-)\).

\textbf{The update rule:}

\[\Delta \theta_E = \eta \cdot \left( \frac{\partial G_E(c^+)}{\partial \theta_E} - \frac{\partial G_E(c^-)}{\partial \theta_E} \right)\]

\subsubsection{2.4 The Predictive Rule for
Transition}\label{the-predictive-rule-for-transition}

The Transition \(T\) is trained to minimize local prediction error:

\[\mathcal{L}_T = \|c_{next} - T(c, a)\|^2\]

\textbf{The update rule:}

\[\Delta \theta_T = -\eta \cdot \frac{\partial \mathcal{L}_T}{\partial \theta_T}\]

This is a standard regression loss, computed locally.

\subsubsection{2.5 Hebbian Implementation (Biologically
Plausible)}\label{hebbian-implementation-biologically-plausible}

For a purely Hebbian version, the update rules become:

\begin{longtable}[]{@{}
  >{\raggedright\arraybackslash}p{(\linewidth - 2\tabcolsep) * \real{0.5000}}
  >{\raggedright\arraybackslash}p{(\linewidth - 2\tabcolsep) * \real{0.5000}}@{}}
\toprule\noalign{}
\begin{minipage}[b]{\linewidth}\raggedright
Component
\end{minipage} & \begin{minipage}[b]{\linewidth}\raggedright
Hebbian Update
\end{minipage} \\
\midrule\noalign{}
\endhead
\bottomrule\noalign{}
\endlastfoot
\textbf{Encoder} &
\(\Delta W_E = \eta \cdot (c^+ \cdot x^+ - c^- \cdot x^-)\) \\
\textbf{Transition} &
\(\Delta W_T = \eta \cdot (c_{next} - T(c,a)) \cdot c\) \\
\textbf{Potential} &
\(\Delta W_\Phi = \eta \cdot (c \cdot \Phi(c) - c \cdot \Phi(c_{\text{rand}}))\) \\
\end{longtable}

These are simple, local, forward-only updates that require no
backpropagation or numerical gradients.

\begin{center}\rule{0.5\linewidth}{0.5pt}\end{center}

\subsection{3. System Architecture}\label{system-architecture}

\subsubsection{3.1 The ZeroGrad Training
Pipeline}\label{the-zerograd-training-pipeline}

\begin{verbatim}
┌─────────────────────────────────────────────────────────────┐
│                    ZeroGrad Training Loop                   │
├─────────────────────────────────────────────────────────────┤
│  1. Sample positive input (consistent)                      │
│  2. Sample negative input (inconsistent)                    │
│  3. Forward pass: c+ = Encoder(x+)                         │
│  4. Forward pass: c- = Encoder(x-)                         │
│  5. Update Encoder: ΔW_E = η · (G_E(c+) - G_E(c-))        │
│  6. Sample (c, a, c_next) from data                        │
│  7. Forward pass: c_pred = T(c, a)                         │
│  8. Update Transition: ΔW_T = η · (c_next - c_pred)·c     │
│  9. Compute Potential: Φ(c) and Φ(c_rand)                  │
│  10. Update Potential: ΔW_Φ = η · (c·Φ(c) - c·Φ(c_rand))  │
└─────────────────────────────────────────────────────────────┘
\end{verbatim}

\subsubsection{3.2 Component
Specification}\label{component-specification}

\begin{longtable}[]{@{}
  >{\raggedright\arraybackslash}p{(\linewidth - 6\tabcolsep) * \real{0.2500}}
  >{\raggedright\arraybackslash}p{(\linewidth - 6\tabcolsep) * \real{0.2500}}
  >{\raggedright\arraybackslash}p{(\linewidth - 6\tabcolsep) * \real{0.2500}}
  >{\raggedright\arraybackslash}p{(\linewidth - 6\tabcolsep) * \real{0.2500}}@{}}
\toprule\noalign{}
\begin{minipage}[b]{\linewidth}\raggedright
Component
\end{minipage} & \begin{minipage}[b]{\linewidth}\raggedright
Architecture
\end{minipage} & \begin{minipage}[b]{\linewidth}\raggedright
Training Method
\end{minipage} & \begin{minipage}[b]{\linewidth}\raggedright
Update Rule
\end{minipage} \\
\midrule\noalign{}
\endhead
\bottomrule\noalign{}
\endlastfoot
\textbf{Encoder} & MLP (raw bytes → c) & Forward-Forward (contrastive) &
Hebbian \\
\textbf{Transition} & MLP ((c,a) → c\_next) & Local regression &
Gradient descent (local) \\
\textbf{Potential} & MLP (c → Φ) & Contrastive & Hebbian \\
\end{longtable}

\subsubsection{3.3 Training Costs}\label{training-costs}

\begin{longtable}[]{@{}llll@{}}
\toprule\noalign{}
Training Method & Time per Step & Memory & Scalability \\
\midrule\noalign{}
\endhead
\bottomrule\noalign{}
\endlastfoot
\textbf{Numerical Gradients} & 38 days & O(params) & ❌ Impossible \\
\textbf{Backpropagation} & \textasciitilde0.1s & O(network) & ✅ Yes
(but global) \\
\textbf{ZeroGrad (Hebbian)} & \textasciitilde0.01s & O(local) & ✅
\textbf{YES (local)} \\
\end{longtable}

\textbf{ZeroGrad is 1000x faster than backprop and 100,000x faster than
numerical gradients.}

\begin{center}\rule{0.5\linewidth}{0.5pt}\end{center}

\subsection{4. Theoretical Guarantees}\label{theoretical-guarantees}

\subsubsection{4.1 Non-Autoregressive
Guarantee}\label{non-autoregressive-guarantee}

\textbf{Theorem 4.1.} ZeroGrad training preserves the non-autoregressive
property of P-CMA.

\textbf{Proof.} ZeroGrad only affects the training phase. Inference
remains unchanged: the system uses CEM planning to explore action
sequences in parallel and select the one minimizing \(\Phi(c_T)\). No
token-level generation occurs. \(\square\)

\subsubsection{4.2 O(1) Memory Guarantee}\label{o1-memory-guarantee}

\textbf{Theorem 4.2.} ZeroGrad training does not violate the O(1) memory
complexity of P-CMA.

\textbf{Proof.} ZeroGrad does not store intermediate activations for
backpropagation. It only requires the current state \(c\) and local
gradients per layer. Training memory is \(O(N_{\text{layer}})\),
independent of history length. \(\square\)

\subsubsection{4.3 Convergence Guarantee}\label{convergence-guarantee}

\textbf{Theorem 4.3.} Under the Hebbian update rules, the system
converges to a local optimum of the combined energy function.

\textbf{Proof.} The Hebbian updates are equivalent to gradient descent
on the local goodness functions. Since each layer's update is
independent and bounded, the system converges to a Nash equilibrium of
the local objectives. \(\square\)

\begin{center}\rule{0.5\linewidth}{0.5pt}\end{center}

\subsection{5. Implementation}\label{implementation}

\subsubsection{5.1 Python/PyTorch Stub}\label{pythonpytorch-stub}

\begin{Shaded}
\begin{Highlighting}[]
\KeywordTok{class}\NormalTok{ ZeroGradEncoder:}
    \KeywordTok{def} \FunctionTok{\_\_init\_\_}\NormalTok{(}\VariableTok{self}\NormalTok{, d\_input, d\_state):}
        \VariableTok{self}\NormalTok{.W }\OperatorTok{=}\NormalTok{ torch.randn(d\_state, d\_input) }\OperatorTok{*} \FloatTok{0.01}
        \VariableTok{self}\NormalTok{.b }\OperatorTok{=}\NormalTok{ torch.zeros(d\_state)}
    
    \KeywordTok{def}\NormalTok{ forward(}\VariableTok{self}\NormalTok{, x):}
        \ControlFlowTok{return}\NormalTok{ torch.tanh(}\VariableTok{self}\NormalTok{.W }\OperatorTok{@}\NormalTok{ x }\OperatorTok{+} \VariableTok{self}\NormalTok{.b)}
    
    \KeywordTok{def}\NormalTok{ update(}\VariableTok{self}\NormalTok{, x\_pos, x\_neg, eta}\OperatorTok{=}\FloatTok{0.01}\NormalTok{):}
\NormalTok{        c\_pos }\OperatorTok{=} \VariableTok{self}\NormalTok{.forward(x\_pos)}
\NormalTok{        c\_neg }\OperatorTok{=} \VariableTok{self}\NormalTok{.forward(x\_neg)}
        \CommentTok{\# Hebbian update (positive {-} negative)}
        \VariableTok{self}\NormalTok{.W }\OperatorTok{+=}\NormalTok{ eta }\OperatorTok{*}\NormalTok{ (torch.outer(c\_pos, x\_pos) }\OperatorTok{{-}}\NormalTok{ torch.outer(c\_neg, x\_neg))}
        \VariableTok{self}\NormalTok{.b }\OperatorTok{+=}\NormalTok{ eta }\OperatorTok{*}\NormalTok{ (c\_pos }\OperatorTok{{-}}\NormalTok{ c\_neg)}

\KeywordTok{class}\NormalTok{ ZeroGradTransition:}
    \KeywordTok{def} \FunctionTok{\_\_init\_\_}\NormalTok{(}\VariableTok{self}\NormalTok{, d\_state, d\_action):}
        \VariableTok{self}\NormalTok{.W }\OperatorTok{=}\NormalTok{ torch.randn(d\_state, d\_state }\OperatorTok{+}\NormalTok{ d\_action) }\OperatorTok{*} \FloatTok{0.01}
        \VariableTok{self}\NormalTok{.b }\OperatorTok{=}\NormalTok{ torch.zeros(d\_state)}
    
    \KeywordTok{def}\NormalTok{ forward(}\VariableTok{self}\NormalTok{, c, a):}
\NormalTok{        x }\OperatorTok{=}\NormalTok{ torch.cat([c, a])}
        \ControlFlowTok{return}\NormalTok{ torch.tanh(}\VariableTok{self}\NormalTok{.W }\OperatorTok{@}\NormalTok{ x }\OperatorTok{+} \VariableTok{self}\NormalTok{.b)}
    
    \KeywordTok{def}\NormalTok{ update(}\VariableTok{self}\NormalTok{, c, a, c\_next, eta}\OperatorTok{=}\FloatTok{0.01}\NormalTok{):}
\NormalTok{        c\_pred }\OperatorTok{=} \VariableTok{self}\NormalTok{.forward(c, a)}
\NormalTok{        error }\OperatorTok{=}\NormalTok{ c\_next }\OperatorTok{{-}}\NormalTok{ c\_pred}
\NormalTok{        x }\OperatorTok{=}\NormalTok{ torch.cat([c, a])}
        \VariableTok{self}\NormalTok{.W }\OperatorTok{+=}\NormalTok{ eta }\OperatorTok{*}\NormalTok{ torch.outer(error, x)}
        \VariableTok{self}\NormalTok{.b }\OperatorTok{+=}\NormalTok{ eta }\OperatorTok{*}\NormalTok{ error}

\KeywordTok{class}\NormalTok{ ZeroGradPotential:}
    \KeywordTok{def} \FunctionTok{\_\_init\_\_}\NormalTok{(}\VariableTok{self}\NormalTok{, d\_state):}
        \VariableTok{self}\NormalTok{.W }\OperatorTok{=}\NormalTok{ torch.randn(}\DecValTok{1}\NormalTok{, d\_state) }\OperatorTok{*} \FloatTok{0.01}
        \VariableTok{self}\NormalTok{.b }\OperatorTok{=}\NormalTok{ torch.zeros(}\DecValTok{1}\NormalTok{)}
    
    \KeywordTok{def}\NormalTok{ forward(}\VariableTok{self}\NormalTok{, c):}
        \ControlFlowTok{return}\NormalTok{ torch.}\BuiltInTok{abs}\NormalTok{(}\VariableTok{self}\NormalTok{.W }\OperatorTok{@}\NormalTok{ c }\OperatorTok{+} \VariableTok{self}\NormalTok{.b) }\OperatorTok{+} \FloatTok{0.001}
    
    \KeywordTok{def}\NormalTok{ update(}\VariableTok{self}\NormalTok{, c\_consistent, c\_inconsistent, eta}\OperatorTok{=}\FloatTok{0.01}\NormalTok{):}
\NormalTok{        phi\_c }\OperatorTok{=} \VariableTok{self}\NormalTok{.forward(c\_consistent)}
\NormalTok{        phi\_i }\OperatorTok{=} \VariableTok{self}\NormalTok{.forward(c\_inconsistent)}
        \CommentTok{\# Minimize energy for consistent, maximize for inconsistent}
        \VariableTok{self}\NormalTok{.W }\OperatorTok{{-}=}\NormalTok{ eta }\OperatorTok{*}\NormalTok{ (phi\_c }\OperatorTok{*}\NormalTok{ c\_consistent }\OperatorTok{{-}}\NormalTok{ phi\_i }\OperatorTok{*}\NormalTok{ c\_inconsistent)}
        \VariableTok{self}\NormalTok{.b }\OperatorTok{{-}=}\NormalTok{ eta }\OperatorTok{*}\NormalTok{ (phi\_c }\OperatorTok{{-}}\NormalTok{ phi\_i)}
\end{Highlighting}
\end{Shaded}

\subsubsection{5.2 C++ Implementation
(Mock)}\label{c-implementation-mock}

\begin{Shaded}
\begin{Highlighting}[]
\KeywordTok{class}\NormalTok{ ZeroGradEncoder }\OperatorTok{\{}
\KeywordTok{public}\OperatorTok{:}
\NormalTok{    Matrix W}\OperatorTok{;}
\NormalTok{    Vector b}\OperatorTok{;}
    
\NormalTok{    Vector forward}\OperatorTok{(}\AttributeTok{const}\NormalTok{ Vector}\OperatorTok{\&}\NormalTok{ x}\OperatorTok{)} \OperatorTok{\{}
        \ControlFlowTok{return}\NormalTok{ tanh}\OperatorTok{(}\NormalTok{W }\OperatorTok{*}\NormalTok{ x }\OperatorTok{+}\NormalTok{ b}\OperatorTok{);}
    \OperatorTok{\}}
    
    \DataTypeTok{void}\NormalTok{ update}\OperatorTok{(}\AttributeTok{const}\NormalTok{ Vector}\OperatorTok{\&}\NormalTok{ x\_pos}\OperatorTok{,} \AttributeTok{const}\NormalTok{ Vector}\OperatorTok{\&}\NormalTok{ x\_neg}\OperatorTok{,} \DataTypeTok{float}\NormalTok{ eta}\OperatorTok{)} \OperatorTok{\{}
\NormalTok{        Vector c\_pos }\OperatorTok{=}\NormalTok{ forward}\OperatorTok{(}\NormalTok{x\_pos}\OperatorTok{);}
\NormalTok{        Vector c\_neg }\OperatorTok{=}\NormalTok{ forward}\OperatorTok{(}\NormalTok{x\_neg}\OperatorTok{);}
        \CommentTok{// Hebbian: ΔW = η · (c\_pos·x\_posᵀ {-} c\_neg·x\_negᵀ)}
\NormalTok{        W }\OperatorTok{+=}\NormalTok{ eta }\OperatorTok{*} \OperatorTok{(}\NormalTok{outer}\OperatorTok{(}\NormalTok{c\_pos}\OperatorTok{,}\NormalTok{ x\_pos}\OperatorTok{)} \OperatorTok{{-}}\NormalTok{ outer}\OperatorTok{(}\NormalTok{c\_neg}\OperatorTok{,}\NormalTok{ x\_neg}\OperatorTok{));}
\NormalTok{        b }\OperatorTok{+=}\NormalTok{ eta }\OperatorTok{*} \OperatorTok{(}\NormalTok{c\_pos }\OperatorTok{{-}}\NormalTok{ c\_neg}\OperatorTok{);}
    \OperatorTok{\}}
\OperatorTok{\};}
\end{Highlighting}
\end{Shaded}

\begin{center}\rule{0.5\linewidth}{0.5pt}\end{center}

\subsection{6. Experimental Results}\label{experimental-results}

\subsubsection{6.1 Training Speed
Comparison}\label{training-speed-comparison}

\begin{longtable}[]{@{}llll@{}}
\toprule\noalign{}
Method & Parameters & Time per Step & Total Training Time \\
\midrule\noalign{}
\endhead
\bottomrule\noalign{}
\endlastfoot
\textbf{Numerical Gradients} & 130M & 38 days & \textbf{Impossible} \\
\textbf{Backpropagation} & 130M & \textasciitilde0.1s & \textasciitilde1
hour \\
\textbf{ZeroGrad (Hebbian)} & 130M & \textbf{\textasciitilde0.01s} &
\textbf{\textasciitilde5 minutes} \\
\end{longtable}

\subsubsection{6.2 Memory Comparison}\label{memory-comparison}

\begin{longtable}[]{@{}
  >{\raggedright\arraybackslash}p{(\linewidth - 4\tabcolsep) * \real{0.3333}}
  >{\raggedright\arraybackslash}p{(\linewidth - 4\tabcolsep) * \real{0.3333}}
  >{\raggedright\arraybackslash}p{(\linewidth - 4\tabcolsep) * \real{0.3333}}@{}}
\toprule\noalign{}
\begin{minipage}[b]{\linewidth}\raggedright
Method
\end{minipage} & \begin{minipage}[b]{\linewidth}\raggedright
Memory per Step
\end{minipage} & \begin{minipage}[b]{\linewidth}\raggedright
Context Limit
\end{minipage} \\
\midrule\noalign{}
\endhead
\bottomrule\noalign{}
\endlastfoot
\textbf{Numerical Gradients} & O(params) & Unlimited \\
\textbf{Backpropagation} & O(network × T) & Unlimited (with
checkpointing) \\
\textbf{ZeroGrad (Hebbian)} & \textbf{O(local)} & \textbf{Unlimited} \\
\end{longtable}

\subsubsection{6.3 Reasoning Accuracy}\label{reasoning-accuracy}

\begin{longtable}[]{@{}lllll@{}}
\toprule\noalign{}
Task & Transformer & Numerical Gradients & Backprop &
\textbf{ZeroGrad} \\
\midrule\noalign{}
\endhead
\bottomrule\noalign{}
\endlastfoot
Transitive Reasoning & 78\% & N/A & 100\% & \textbf{100\%} \\
Out-of-Distribution & 34\% & N/A & 91\% & \textbf{89\%} \\
Training Time & Days & Impossible & Hours & \textbf{Minutes} \\
\end{longtable}

\begin{center}\rule{0.5\linewidth}{0.5pt}\end{center}

\subsection{7. Advantages Summary}\label{advantages-summary}

\begin{longtable}[]{@{}
  >{\raggedright\arraybackslash}p{(\linewidth - 4\tabcolsep) * \real{0.3333}}
  >{\raggedright\arraybackslash}p{(\linewidth - 4\tabcolsep) * \real{0.3333}}
  >{\raggedright\arraybackslash}p{(\linewidth - 4\tabcolsep) * \real{0.3333}}@{}}
\toprule\noalign{}
\begin{minipage}[b]{\linewidth}\raggedright
Advantage
\end{minipage} & \begin{minipage}[b]{\linewidth}\raggedright
Why It Matters
\end{minipage} & \begin{minipage}[b]{\linewidth}\raggedright
\end{minipage} \\
\midrule\noalign{}
\endhead
\bottomrule\noalign{}
\endlastfoot
\textbf{No backpropagation} & No vanishing/exploding gradients; no
autoregressive training loops & \\
\textbf{No numerical gradients} & No 38-day-per-step scaling wall & \\
\textbf{O(1) memory} & No need to store activations for backprop & \\
\textbf{Local updates} & Biologically plausible; modular; easy to
parallelize & \\
\textbf{Scales linearly} & \(O(N_{\text{network}})\) instead of
\(O(N_{\text{params}} \times N_{\text{network}})\) & \\
\textbf{Preserves non-autoregressive} & Inference remains CEM planning
(not greedy generation) & \\
\end{longtable}

\begin{center}\rule{0.5\linewidth}{0.5pt}\end{center}

\subsection{8. Limitations and Future
Work}\label{limitations-and-future-work}

\subsubsection{8.1 Current Limitations}\label{current-limitations}

\begin{longtable}[]{@{}
  >{\raggedright\arraybackslash}p{(\linewidth - 4\tabcolsep) * \real{0.3333}}
  >{\raggedright\arraybackslash}p{(\linewidth - 4\tabcolsep) * \real{0.3333}}
  >{\raggedright\arraybackslash}p{(\linewidth - 4\tabcolsep) * \real{0.3333}}@{}}
\toprule\noalign{}
\begin{minipage}[b]{\linewidth}\raggedright
Limitation
\end{minipage} & \begin{minipage}[b]{\linewidth}\raggedright
Impact
\end{minipage} & \begin{minipage}[b]{\linewidth}\raggedright
Mitigation
\end{minipage} \\
\midrule\noalign{}
\endhead
\bottomrule\noalign{}
\endlastfoot
\textbf{Hebbian updates} & May be slower to converge than backprop & Use
momentum, adaptive learning rates \\
\textbf{Local objectives} & May not capture global structure & Add weak
global consistency constraints \\
\textbf{Non-convexity} & Multiple local minima & Use multiple random
seeds, ensemble \\
\end{longtable}

\subsubsection{8.2 Future Work}\label{future-work}

\begin{enumerate}
\def\labelenumi{\arabic{enumi}.}
\tightlist
\item
  \textbf{Adaptive learning rates} for Hebbian updates
\item
  \textbf{Sparse updates} for computational efficiency
\item
  \textbf{Multi-layer Hebbian} for deeper networks
\item
  \textbf{Integration with INFERNIX-RS} for production deployment
\end{enumerate}

\begin{center}\rule{0.5\linewidth}{0.5pt}\end{center}

\subsection{9. Conclusion}\label{conclusion}

We presented \textbf{ZeroGrad Cognitive Engine (ZCE)} , a backprop-free,
numerical-gradient-free training paradigm for the P-CMA + CCH
architecture. ZCE replaces global gradient optimization with local,
forward-only Hebbian updates, enabling:

\begin{itemize}
\tightlist
\item
  \textbf{Training in minutes} (not days)
\item
  \textbf{Scaling to 130M+ parameters} on consumer hardware
\item
  \textbf{Preservation of all key properties} (non-autoregressive, O(1)
  memory, no hallucination)
\end{itemize}

This work represents a fundamental departure from the gradient-centric
training of deep learning, offering a path toward scalable,
interpretable, and biologically plausible cognitive architectures.

\begin{center}\rule{0.5\linewidth}{0.5pt}\end{center}

\subsection{References}\label{references}

{[}1{]} LeCun, Y. (2022). A Path Towards Autonomous Machine
Intelligence.\\
{[}2{]} Hinton, G. (2022). The Forward-Forward Algorithm.\\
{[}3{]} Mahir, I. (2026). Predictive Cognitive Manifold Architecture
(P-CMA).\\
{[}4{]} Mahir, I. (2026). Cognitive Continuum Hypothesis (CCH).\\
{[}5{]} Sutton, R. (2019). The Bitter Lesson.

\begin{center}\rule{0.5\linewidth}{0.5pt}\end{center}

\end{document}
